\documentclass[12pt]{article}
\usepackage[T1]{fontenc}
\usepackage[utf8]{inputenc}
\usepackage{lmodern}
\usepackage{textcomp}
\usepackage{enumitem}
\usepackage{geometry}
\usepackage{graphicx}
\usepackage{amsmath, amssymb}
\geometry{margin=1in}
\begin{document}
\begin{center}
Health Sciences - Graduate Level Assessment 12 \
Total Marks: 40 \
Answer all questions.
\end{center}

\section*{Multiple Choice (10 marks)}
\begin{enumerate}[label=\arabic*.]
\item Which of the following confirmed values meet the diagnostic threshold for diabetes?
\begin{enumerate}[label=(\alph*)]
    \item fasting blood glucose \ensuremath{\geq} 160 mg/dl
    \item fasting blood glucose ? 140 mg/dl
    \item fasting blood glucose \ensuremath{\geq} 126 mg/dl
    \item random glucose \textgreater{} 180 mg/dl
    \item fasting blood glucose ? 120 mg/dl
\end{enumerate}
\item Simple tandem repeat polymorphisms in humans are most useful for
\begin{enumerate}[label=(\alph*)]
    \item studying the effects of radiation exposure
    \item identifying the presence of bacterial infections
    \item solving criminal and paternity cases
    \item determining blood type
    \item transferring disease resistance factors into bone marrow cells
\end{enumerate}
\item Proprioceptive nerve endings in synovial joints are located in
\begin{enumerate}[label=(\alph*)]
    \item synovial fluid and articular discs.
    \item capsule and ligaments.
    \item articular cartilage and ligaments.
    \item capsule and articular cartilage.
    \item synovial membrane and capsule.
\end{enumerate}
\item Severe acute malnutrition in young children is defined as:
\begin{enumerate}[label=(\alph*)]
    \item Weight-for-age Z score \textless{}-2 and oedema
    \item Height-for-age Z score \textless{}-2 or weight-for-height Z score \textless{}-2 and oedema
    \item Height-for-age Z score \textless{}-3 and oedema
    \item Weight-for-age Z score \textless{}-2 or height-for-age Z score \textless{}-2 or oedema
    \item Height-for-age Z score \textless{}-3 or weight-for-age Z score \textless{}-3 and oedema
\end{enumerate}
\item What is the difference between a male and a female catheter?
\begin{enumerate}[label=(\alph*)]
    \item Female catheters are used more frequently than male catheters.
    \item Male catheters are bigger than female catheters.
    \item Male catheters are more flexible than female catheters.
    \item Male catheters are made from a different material than female catheters.
    \item Female catheters are longer than male catheters.
\end{enumerate}
\end{enumerate}

\section*{True / False (10 marks)}
\textit{Answer True or False}
\begin{enumerate}[label=\arabic*.]
\setcounter{enumi}{5}
\item True or False: Grade 2/5 power represents movement of the arm only against resistance.
\item True or False: Having a lower level of education is associated with a decreased risk of Alzheimer's disease.
\item True or False: The maxillary bone has a process that connects to the zygomatic bone, forming the zygomatic arch.
\item True or False: Type I muscle fibers are characterized as red, oxidative, and slow contracting.
\item True or False: The concentration of a food substance is a critical factor in evaluating its safety based on toxicity and dose-response data.
\end{enumerate}

\section*{Long Form Response (20 marks)}
\textit{Answer each question with a clear and complete explanation.}
\begin{enumerate}[label=\arabic*.]
\setcounter{enumi}{10}
\item Discuss the genetic basis of DiGeorge/Shprintzen syndrome, including the implications of the chromosomal deletion involved and its effects on physiological and developmental outcomes.
\item Discuss the role of ubiquitin in proteolysis, particularly focusing on its attachment to proteins designated for synthesis versus those targeted for catabolism.
\end{enumerate}

\vfill
\noindent\textit{End of Paper}
\end{document}